\ReportSection{4. Structure of the Story Test Suite}
\ReportSubsection{4.1 Feature Files and Step Definitions}
Cucumber.js reads the feature files and matches their steps to JavaScript
functions in \Path{steps/}. Each story has its own step-definition file. We put
common fixture setup, response checks, and cleanup in \Path{steps/stepDefs.js}
so that several stories can use the same steps. \Path{TestUtil.js} handles HTTP
requests and record lookups, while Chai assertions check the responses and stored
results.

\begin{center}
\begin{tikzpicture}[node distance=10bp]
\node[process,text width=\dimexpr\linewidth-30bp\relax] (features)
{\textbf{Feature files} --- user goals, \textit{Backgrounds}, \textit{Scenario Outlines}, Examples};
\node[process,text width=\dimexpr\linewidth-30bp\relax,below=of features] (steps)
{\textbf{Story steps and shared library} --- fixtures, requests, correctness assertions};
\node[process,text width=\dimexpr\linewidth-30bp\relax,below=of steps] (helpers)
{\textbf{TestUtil.js and setup.js} --- HTTP helpers, snapshots, readiness, restart, cleanup};
\node[terminal,text width=\dimexpr\linewidth-30bp\relax,below=of helpers] (api)
{\textbf{Todo Manager 1.5.5} --- local REST API at port 4567};
\draw[arrow] (features)--(steps);
\draw[arrow] (steps)--(helpers);
\draw[arrow] (helpers)--(api);
\end{tikzpicture}
\captionof{figure}{How acceptance scenarios reach the API and assess its behavior.}
\label{fig:architecture}
\end{center}

\ReportSubsection{4.2 Setup, Assertions, and Cleanup}
The shared readiness step stops any existing API instance and starts the bundled
JAR. It waits for HTTP 200 and a todos array, then records the seeded todos,
projects, and categories. The remaining \textit{Background} steps create the fixtures for
the example. Each scenario keeps its responses and selected IDs in Cucumber's
\Path{this} state, which is separate for each case.

After a write, the steps read back the data to check what was stored. Where a
change should affect only one record, the tests also compare the other records
with the saved state. This allows the tests to detect unwanted side effects
that would not be apparent from the response status alone.

The \textit{After} hook runs even if an assertion fails. It stops and restarts the JAR,
compares the seeded records with the initial snapshot, and stops the service
again. Before comparison, the helpers sort records and relationship references
by ID so that a different array order does not cause a false failure.

\ReportSubsection{4.3 Independence and Code Organization}
Each example gets a fresh application process and creates its own fixtures.
Helpers use returned IDs or look up records by title, so the tests do not need
to predict which ID will be allocated next. Shared request and payload functions
keep the HTTP handling in one place. Conversion of Gherkin strings to booleans
is limited to the boolean fields, leaving invalid values intact for error tests.
HTTP requests have a five-second timeout, with separate time limits for readiness
and application startup and shutdown.
